% !TeX root = ../Book1.tex
\begin{tikzpicture}[scale=.94, every node/.style={font=\small}]
	\begin{scope}[shift={(-4.4,0)}]
		\fill[blue!9] (0,0) -- (0:1.17) arc (0:120:1.17) -- cycle;
		\draw (0,0) circle (1.17);
		\draw[blue!55] (0,0) -- (0:1.17) (0,0) -- (120:1.17) (0,0) -- (240:1.17);
		\foreach \a in {30,150,270} {\fill[blue!70!black] (\a:.75) circle (1.8pt);}
		\fill (0,0) circle (1.8pt);
		\node[below left] at (0,0) {\(0\)};
		\node at (0,-1.55) {圆盘中的旋转轨道};
	\end{scope}
	\draw[->,thick] (-2.95,0) -- (-1.45,0) node[midway,above] {\(z\mapsto z^m\)};
	\begin{scope}
		\fill[blue!5] (0,0) circle (1.17);
		\draw (0,0) circle (1.17);
		\fill (0,0) circle (1.8pt);
		\node[below] at (0,-.05) {\(\mathbb Z/m\mathbb Z\)};
		\fill (.55,.58) circle (1.8pt);
		\node[above left] at (.55,.58) {\(\{e\}\)};
		\node at (0,-1.55) {旋转商：仍是圆盘};
	\end{scope}
	\begin{scope}[shift={(4.3,-.35)}]
		\fill[blue!5] (-1.17,0) arc (180:0:1.17) -- cycle;
		\draw (-1.17,0) arc (180:0:1.17) -- cycle;
		\draw[blue!65,thick] (-1.17,0) -- (1.17,0);
		\fill (0,0) circle (1.8pt);
		\node[below] at (0,0) {\(\mathbb Z/2\mathbb Z\)};
		\fill (.2,.65) circle (1.8pt);
		\node[above left] at (.2,.65) {\(\{e\}\)};
		\node at (0,-1.2) {反射商：半圆盘};
	\end{scope}
\end{tikzpicture}
